% ==============================================================
%  Naawbi LMS --- CS3009 Software Engineering (Spring 2026)
%  Deliverable 3 — Sprint 3 + Deliverable Document  |  Team Ibwaan
% ==============================================================
\documentclass[12pt, a4paper]{article}

% ── Packages ──────────────────────────────────────────────────
\usepackage[T1]{fontenc}
\usepackage[utf8]{inputenc}
\usepackage{lmodern}
\usepackage[top=2.3cm, bottom=2.3cm, left=2.3cm, right=2.3cm, headheight=28pt]{geometry}
\usepackage{microtype}
\usepackage{setspace}
\usepackage{parskip}
\usepackage{titlesec}
\usepackage{titling}
\usepackage{fancyhdr}
\usepackage{graphicx}
\usepackage{booktabs}
\usepackage{longtable}
% array + tabularx removed: MiKTeX bug with p{} + colortbl from [table]xcolor
\usepackage{multicol}
\usepackage{enumitem}
\usepackage{mdframed}
\usepackage{xcolor}
\usepackage{hyperref}
\usepackage{marvosym}
\newcommand{\faGithub}{\raisebox{-0.1ex}{\Keyboard}}
\usepackage{soul}
\usepackage{float}
\usepackage{caption}
\usepackage{subcaption}
\usepackage[most]{tcolorbox}

% ── Colours ───────────────────────────────────────────────────
\definecolor{primary}{HTML}{1A237E}
\definecolor{accent}{HTML}{283593}
\definecolor{light}{HTML}{E8EAF6}
\definecolor{rule}{HTML}{5C6BC0}
\definecolor{muted}{HTML}{546E7A}
\definecolor{rowalt}{HTML}{F3F4FA}
\definecolor{boxbg}{HTML}{EEF0FB}
\definecolor{success}{HTML}{1B5E20}
\definecolor{warn}{HTML}{B26A00}
\definecolor{danger}{HTML}{B71C1C}

% ── Hyperref ─────────────────────────────────────────────────
\hypersetup{
    colorlinks   = true,
    linkcolor    = primary,
    urlcolor     = accent,
    citecolor    = primary,
    pdftitle     = {Naawbi LMS — Deliverable 3},
    pdfauthor    = {Team Ibwaan},
    pdfsubject   = {CS3009 Software Engineering Spring 2026},
    pdfkeywords  = {LMS, Software Engineering, Naawbi, Ibwaan, Sprint 3},
}

% ── Sections ─────────────────────────────────────────────────
\titleformat{\section}
  {\Large\bfseries\color{primary}}
  {\thesection}{1em}{}
  [\color{rule}\titlerule]

\titleformat{\subsection}
  {\large\bfseries\color{accent}}
  {\thesubsection}{1em}{}

\titleformat{\subsubsection}
  {\normalsize\bfseries\color{accent}}
  {\thesubsubsection}{1em}{}

\titlespacing*{\section}{0pt}{1.6ex plus .5ex minus .2ex}{0.8ex}
\titlespacing*{\subsection}{0pt}{1.0ex plus .3ex minus .1ex}{0.4ex}

% ── Header / Footer ───────────────────────────────────────────
\pagestyle{fancy}
\fancyhf{}
\renewcommand{\headrulewidth}{0.4pt}
\renewcommand{\headrule}{\hbox to\headwidth{\color{rule}\leaders\hrule height \headrulewidth\hfill}}
\fancyhead[L]{\small\color{muted}\textbf{Naawbi LMS} — Deliverable 3}
\fancyhead[R]{\small\color{muted}CS3009 Spring 2026}
\fancyfoot[C]{\small\color{muted}\thepage}

% ── Lists ─────────────────────────────────────────────────────
\setlist[itemize]{leftmargin=1.5em, itemsep=2pt, topsep=4pt}
\setlist[enumerate]{leftmargin=1.5em, itemsep=2pt, topsep=4pt}

% ── Image fallback ────────────────────────────────────────────
\newcommand{\imgOrPlaceholder}[3][width=0.85\textwidth]{%
  \IfFileExists{#2}{%
    \includegraphics[#1]{#2}%
  }{%
    \begin{center}%
    \fboxsep=8pt\fbox{\parbox{0.7\textwidth}{\centering\color{muted}%
      \textit{[Placeholder — image not yet available]}\\[4pt]%
      \footnotesize #3}}%
    \end{center}%
  }%
}

% ── Custom environments ───────────────────────────────────────
\tcbuselibrary{skins, breakable}

\newtcolorbox{infobox}[1][]{
  enhanced, breakable,
  colback=boxbg, colframe=rule,
  arc=4pt, boxrule=0.8pt,
  leftrule=4pt,
  title=#1,
  fonttitle=\bfseries\color{primary},
  #1
}

\newtcolorbox{userstory}[1][]{
  enhanced, breakable,
  colback=white, colframe=rule!40,
  arc=3pt, boxrule=0.5pt,
  top=4pt, bottom=4pt, left=8pt, right=8pt,
  #1
}

% Row-striping for tables
% altrow stripped — rowcolor conflicts with tabularx X columns

\newcommand{\fcol}[1]{\textbf{\color{primary}{#1}}}
\newcommand{\fcolplain}[1]{\textbf{#1}}
\begin{document}

% ═══════════════════════════════════════════════════════════════
%  TITLE PAGE
% ═══════════════════════════════════════════════════════════════
\begin{titlepage}
  \centering

  \vspace*{2.2cm}

  {\small\scshape\color{muted} National University of Computer and Emerging Sciences}\\[0.4em]
  {\small\color{muted} CS3009 --- Software Engineering \quad|\quad Spring 2026}

  \vspace{1.2cm}
  {\color{rule}\hrule height 1.5pt}
  \vspace{1.2cm}

  {\fontsize{56}{62}\selectfont\bfseries\color{primary} Naawbi}\\[0.6em]
  {\Large\color{muted} Learning Management System}

  \vspace{1.0cm}
  {\color{rule!40}\rule{5cm}{0.6pt}}
  \vspace{1.0cm}

  {\LARGE\bfseries\color{primary} Deliverable 3}\\[0.4em]
  {\normalsize\color{muted} Sprint 3 (Final Sprint) \,\textbar\, Scrum Report + Deliverable Document}

  \vspace{2cm}

  \begin{tabular}{r@{\quad}l}
    {\color{muted}\small Team Name}       & {\bfseries Ibwaan} \\[5pt]
    {\color{muted}\small Team Lead}       & {\bfseries Ibraheem Farooq} \\[5pt]
    {\color{muted}\small Members}         & 23i-0708 \textbar\ Ali Aan Khowaja \\
                                          & 23i-0816 \textbar\ Ibraheem Farooq \\
                                          & 23i-0819 \textbar\ Wajih-Ur-Raza Asif \\[5pt]
    {\color{muted}\small Submission Date} & 20 April 2026 \\[5pt]
    {\color{muted}\small Repository}      & \href{https://github.com/aliaankhowaja/Naawbi}{github.com/aliaankhowaja/Naawbi} \\
  \end{tabular}

  \vfill
  {\color{rule!40}\hrule height 0.5pt}
  \vspace{0.6cm}
  {\footnotesize\color{muted}
    \href{https://github.com/ibbi1020}{ibbi1020}\enspace\textbullet\enspace
    \href{https://github.com/wajih-rathore}{wajih-rathore}\enspace\textbullet\enspace
    \href{https://github.com/aliaankhowaja}{aliaankhowaja}
  }
  \vspace{0.6cm}
\end{titlepage}

% ── Table of Contents ─────────────────────────────────────────
\tableofcontents
\newpage

% ═══════════════════════════════════════════════════════════════
%                 PART B — SCRUM (SPRINT 3)
% ═══════════════════════════════════════════════════════════════
\begin{center}
  {\LARGE\bfseries\color{primary} PART B\,\textemdash\,Development Methodology: Scrum}\\[4pt]
  {\color{muted}\small Sprint 3 — the final sprint}\\[4pt]
  {\color{rule!40}\rule{7cm}{0.6pt}}
\end{center}

% ───────────────────────────────────────────────────────────────
\section{Sprint 3 Overview}

\begin{center}
\begin{tabular}{@{}p{0.22\linewidth}@{\hspace{6pt}}p{0.72\linewidth}@{}}
  \toprule
  \fcol{Field} & Value \\
  \midrule
  \fcol{Sprint Duration} & 26 March 2026 – 20 April 2026 (4 weeks) \\
\fcol{Sprint Goal} & Close the platform's remaining feature gaps: implement
                           user authentication (register / login / logout), enforce
                           role-based access control across all protected actions,
                           add file and link attachments to announcements, build the
                           Assignments and People tabs in the course home, and deliver
                           the full instructor grading interface with a role-adaptive
                           gradebook (student personal view and instructor class-wide
                           grid). \\
  \fcol{Modules Developed} & Authentication, Role-Based Access Control (RBAC),
                         Announcement Attachments, Assignments Tab, People Tab,
                         Grading Interface, Gradebook (student + instructor). \\
\fcol{Carried from D2} & US\,1–4 (Authentication), US\,12–13 (Announcement Attachments), US\,19–22 (Grading \& Gradebook). \\
  \fcol{Added in Sprint 3} & US\,15–16 (Assignments Tab, People Tab) — implied by the
                         D2 assignment workflow but never formally scoped before. \\
  \bottomrule
\end{tabular}
\end{center}

% ───────────────────────────────────────────────────────────────
\section{Sprint 3 Backlog}

\subsection{Handoff from Sprint 2}

Deliverable~2's final Trello / GitHub Projects snapshot showed the following
user stories still in the product backlog at sprint close: \textbf{US\,1–4}
(Authentication), \textbf{US\,12–13} (Announcement Attachments), and
\textbf{US\,19–22} (Grading \& Gradebook). During Sprint~3 planning the team
additionally scoped \textbf{US\,15–16} (Assignments tab + People tab).

\subsection{Module Selected for Sprint 3}

The Sprint~3 module is best described as \textbf{``Identity, Access, and
Evaluation''} --- a composite of the three functional concerns that remained
unbuilt after Sprint~2:
\begin{enumerate}
  \item \textbf{Identity:} user accounts, sessions, passwords.
  \item \textbf{Access:} role-based access control across all protected flows.
  \item \textbf{Evaluation:} grading interface, gradebook (personal + class-wide),
        plus the supporting tabs (Assignments tab, People tab).
\end{enumerate}

\subsection{User Stories Selected for Sprint 3}

\begin{enumerate}[label=\textbf{US\,\arabic*}]
  \setcounter{enumi}{0}
  \item As a \emph{User}, I want to register with a username, email, password, and role so that I have my own account on the platform.
  \item As a \emph{User}, I want to log in with my email and password so that I can access my personalised course workspace.
  \item As a \emph{User}, I want to be assigned a role (Instructor or Student) at registration so that the system tailors my experience and enforces appropriate permissions.
  \item As a \emph{User}, I want to log out so that my session is terminated and my account is secure.
  \item As the \emph{System}, I want to enforce role-based access control so that instructors and students only see and use features appropriate to their role.
\end{enumerate}

\begin{enumerate}[label=\textbf{US\,\arabic*}, start=12]
  \item As an \emph{Instructor}, I want to attach a file to an announcement so that I can share reference materials directly from the course stream.
  \item As an \emph{Instructor}, I want to attach a URL to an announcement so that I can share links to external resources without requiring a file upload.
\end{enumerate}

\begin{enumerate}[label=\textbf{US\,\arabic*}, start=15]
  \item \textbf{(Assignments Tab)} As a \emph{User}, I want a dedicated Assignments tab in the course home showing all assignments and their status at a glance, so I do not have to navigate away from the course to see my workload.
  \item \textbf{(People Tab)} As a \emph{User}, I want a People tab in the course home showing everyone enrolled in the course with their role badge, so I know who is in my class.
\end{enumerate}

\begin{enumerate}[label=\textbf{US\,\arabic*}, start=19]
  \item As an \emph{Instructor}, I want to open a grading view for an assignment so that I can see all student submissions in one place.
  \item As an \emph{Instructor}, I want to assign a numeric grade and written feedback to each submission so that students receive clear, actionable results.
  \item As a \emph{Student}, I want to see my grades and instructor feedback for each graded assignment so that I know how I performed.
  \item As an \emph{Instructor}, I want a class-wide gradebook grid showing all students and assignments so that I have a complete overview of course performance without opening submissions individually.
\end{enumerate}

\subsection{Sub User Stories and Task Breakdowns}

{\sloppy
\begin{itemize}
  \item \textbf{US\,1--4 Authentication:} tabbed \texttt{LoginView} (Login/Register), SHA-256 hashing in \texttt{PasswordUtil}, \texttt{User.register()}, \texttt{LoginController}, \texttt{Session} singleton, and a logout handler that clears the session and reloads the login screen.
  \item \textbf{US\,5 RBAC:} enrolment-based catalog filtering in \texttt{CourseCatalogController}; role guard before Create Assignment / Post Announcement / Grade Submissions with Access Denied alert; instructor-only buttons conditionally rendered, not merely disabled.
  \item \textbf{US\,12--13 Announcement Attachments:} \texttt{CreateAnnouncementController} extended with optional FileChooser + URL field; \texttt{announcement\_attachments} table; chips rendered in the course stream (file via \texttt{Desktop.getDesktop()}, link in default browser).
  \item \textbf{US\,15--16 Course Home Tabs:} \texttt{Assignment.fetchWithStatusForUser()} derives status via LEFT JOIN (Submitted / Late / Missing / Not Submitted); \texttt{Course.fetchEnrolledUsers()} powers the People tab; role-adaptive cards (status badge for students, Grade button for instructors).
  \item \textbf{US\,19--20 Grading:} \texttt{grade} + \texttt{feedback} columns added to \texttt{submissions}; \texttt{Submission.fetchByAssignmentId} / \texttt{saveGrade}; two-panel \texttt{GradingView.fxml} with numeric validation against \texttt{total\_points} before a parameterised UPDATE.
  \item \textbf{US\,21--22 Gradebook:} \texttt{fetchGradedForStudent} (personal view) and \texttt{fetchGradebookForCourse} (class grid) --- single-query, no N+1; \texttt{GradebookView} / \texttt{GradebookController} branch on role to render a table or a GridPane.
\end{itemize}
\par}

% ───────────────────────────────────────────────────────────────
\subsection{Structured Specifications for Four User Stories}

% ── Spec 1: US 1–4 Authentication
\begin{userstory}
  \textbf{Spec 1 --- US\,1--4: Authentication (Register + Login + Logout)}
\end{userstory}
\begin{center}
\begin{tabular}{@{}p{0.22\linewidth}@{\hspace{6pt}}p{0.72\linewidth}@{}}
  \toprule
  \fcol{Field} & Detail \\
  \midrule
  \fcol{Actor} & Any user (Student or Instructor) \\
\fcol{Inputs} & \emph{Register:} username, email, password, role (Student / Instructor). \emph{Login:} email, password. \\
  \fcol{Processing} & Registration: validate all fields non-empty; hash password with SHA-256; call \texttt{User.register()} which inserts the user and returns an empty \texttt{Optional} on duplicate username/email. Login: hash submitted password, compare against stored hash, populate \texttt{Session} on match. Logout: \texttt{Session.logout()} clears state; navigation returns to \texttt{LoginView}. \\
\fcol{Outputs} & On register/login success: user lands on \texttt{CourseCatalogView} with an active session. On logout: session cleared; user returns to \texttt{LoginView}. \\
  \fcol{Preconditions} & Register: username and email not already in \texttt{users}. Login: account exists with a stored hash. \\
\fcol{Postconditions} & \texttt{Session} holds \texttt{id, username, email, role}. Logout leaves \texttt{Session} empty; no credentials persist. \\
  \fcol{Exceptions} & Wrong password / unknown email $\rightarrow$ ``Incorrect email or password''. Empty fields $\rightarrow$ ``Please fill in all fields''. Duplicate username/email on register $\rightarrow$ ``Username or email is already taken''. \\
  \bottomrule
\end{tabular}
\end{center}

% ── Spec 2: US 5 RBAC
\begin{userstory}
  \textbf{Spec 2 --- US\,5: Role-Based Access Control}
\end{userstory}
\begin{center}
\begin{tabular}{@{}p{0.22\linewidth}@{\hspace{6pt}}p{0.72\linewidth}@{}}
  \toprule
  \fcol{Field} & Detail \\
  \midrule
  \fcol{Actor} & System (enforced at every protected point) \\
\fcol{Inputs} & \texttt{Session.getInstance().getRole()} --- set at login/register. \\
  \fcol{Processing} & On catalog load: students see only courses with a \texttt{course\_enrollments} record; instructors see courses where \texttt{courses.created\_by\,=\,userId}. On any instructor-only action (Post Announcement, Create Assignment, Grade Submissions) the controller checks role and shows Access Denied for students. Instructor-only UI elements are conditionally rendered (not merely disabled). \\
\fcol{Outputs} & Students see only enrolled courses. Instructor-only features are inaccessible to students both visually and via guard logic. \\
  \fcol{Preconditions} & User is logged in; \texttt{Session} is populated. \\
\fcol{Postconditions} & No student can trigger an instructor action even by direct controller call. \\
  \fcol{Exceptions} & No unauthenticated code path exists after Sprint~3; \texttt{LoginView} is always the first screen. \\
  \bottomrule
\end{tabular}
\end{center}

% ── Spec 3: US 12–13 Attachments
\begin{userstory}
  \textbf{Spec 3 --- US\,12--13: Announcement Attachments (File + Link)}
\end{userstory}
\begin{center}
\begin{tabular}{@{}p{0.22\linewidth}@{\hspace{6pt}}p{0.72\linewidth}@{}}
  \toprule
  \fcol{Field} & Detail \\
  \midrule
  \fcol{Actor} & Instructor (attach); all enrolled users (view) \\
\fcol{Inputs} & Optional: file selected via system FileChooser OR URL typed into a text field. Both are optional --- an announcement with no attachment is valid. \\
  \fcol{Processing} & File: store original filename and path in \texttt{announcement\_attachments} with \texttt{is\_link\,=\,false}. URL: store in \texttt{link\_url} with \texttt{is\_link\,=\,true}. Both records are linked by foreign key to the announcement. A single announcement can carry one file and one link attachment independently. \\
\fcol{Outputs} & Attachment chips rendered inline below announcement body. File chip opens the file via \texttt{Desktop.getDesktop()}; link chip opens the URL in the default browser. \\
  \fcol{Preconditions} & Instructor is creating or has created an announcement. \\
\fcol{Postconditions} & Attachment record persisted; visible to all enrolled students on the next stream load. \\
  \fcol{Exceptions} & No file/URL provided $\rightarrow$ announcement saves normally with no attachment record. Cleared attachment fields are ignored. \\
  \bottomrule
\end{tabular}
\end{center}

% ── Spec 4: US 19–20 Grading
\begin{userstory}
  \textbf{Spec 4 --- US\,19--20: Instructor Grading Interface}
\end{userstory}
\begin{center}
\begin{tabular}{@{}p{0.22\linewidth}@{\hspace{6pt}}p{0.72\linewidth}@{}}
  \toprule
  \fcol{Field} & Detail \\
  \midrule
  \fcol{Actor} & Instructor \\
\fcol{Inputs} & Selected submission from the ListView; numeric grade (required); feedback text (optional). \\
  \fcol{Processing} & On open: load all submissions via \texttt{Submission.fetchByAssignmentId()}; pre-fill grade/feedback if a prior grade exists. On Save: validate grade is numeric and within $0$ -- \texttt{total\_points}; call \texttt{Submission.saveGrade(submissionId, grade, feedback)} which issues a parameterised UPDATE. \\
\fcol{Outputs} & \texttt{submissions.grade} and \texttt{submissions.feedback} updated. ListView reflects graded state. Changes immediately visible in student gradebook. \\
  \fcol{Preconditions} & Instructor is logged in and owns the course; assignment exists. \\
\fcol{Postconditions} & Submission row has non-null \texttt{grade} and optionally non-null \texttt{feedback}. \\
  \fcol{Exceptions} & No submissions $\rightarrow$ empty ListView with empty-state label. Non-numeric grade or grade $>$ \texttt{total\_points} $\rightarrow$ inline validation error; Save is blocked. \\
  \bottomrule
\end{tabular}
\end{center}

% ───────────────────────────────────────────────────────────────
\section{Scrum Board (GitHub Projects)}

Sprint~3 was tracked on the team's GitHub Projects board. The screenshots
below are taken directly from the live project board at sprint close and
show the full movement of issues across all three phases of the sprint:
issues that were added during Sprint planning (26~Mar), worked on through
mid-sprint, and finally resolved into the \textbf{Done} column by sprint
completion (20~Apr). No leftover tasks remained in Todo or In~Progress.

\begin{figure}[H]
  \centering
  \imgOrPlaceholder[width=0.95\textwidth]{assets/sprintboard1.png}{GitHub Projects Scrum Board}
  \caption{Sprint~3 scrum board on GitHub Projects --- top view of the kanban columns (\textbf{Todo}, \textbf{In Progress}, \textbf{Done}) with Sprint~3 issues.}
\end{figure}

\begin{figure}[H]
  \centering
  \imgOrPlaceholder[width=0.95\textwidth]{assets/sprintboard2.png}{GitHub Projects Scrum Board (continued)}
  \caption{Sprint~3 scrum board --- remaining issues and completed work in the \textbf{Done} column at sprint close (20~Apr~2026).}
\end{figure}

% ───────────────────────────────────────────────────────────────
\section{Sprint 3 Burn-Down Chart}

\begin{figure}[H]
  \centering
  \imgOrPlaceholder[width=0.95\textwidth]{assets/burndown-chart.png}{Sprint 3 Burn-Down Chart}
  \caption{Sprint~3 burn-down chart exported from the team's GitHub Projects board, showing planned vs.\ actual remaining work across the sprint window (26~Mar~--~20~Apr~2026).}
\end{figure}

% ───────────────────────────────────────────────────────────────
\section{Testing Activities for Sprint 3}
\label{sec:testing}

\subsection{Software Test Plan}

\begin{center}
\begin{tabular}{@{}p{0.22\linewidth}@{\hspace{6pt}}p{0.72\linewidth}@{}}
  \toprule
  \fcol{Field} & Detail \\
  \midrule
  \fcol{Scope} & Sprint~3 features: Authentication (register/login/logout), RBAC, Announcement Attachments, Assignments Tab, People Tab, Grading Interface, Gradebook. \\
\fcol{Objectives} & Verify each feature meets its acceptance criteria; detect regressions in Sprint~1/2 features; validate RBAC enforcement across all protected actions. \\
  \fcol{Test Types} & Black-box functional testing (equivalence partitioning, boundary value analysis, error guessing); white-box structural testing (branch coverage on key controller methods). \\
\fcol{Test Environment} & Windows~11, Java~24, JavaFX~21, PostgreSQL~16 --- seeded with \texttt{data/seed.sql} (3 users: \texttt{ali} = instructor, \texttt{ibbi} + \texttt{sara} = students; 2 courses; 6 assignments; 6 submissions in mixed states). \\
  \fcol{Entry Criteria} & \texttt{make seed} succeeds; \texttt{make compile} produces zero errors; application launches to \texttt{LoginView}. \\
\fcol{Exit Criteria} & All P1 (critical) test cases pass; no unfixed P1 bugs; all identified P2 bugs documented with status. \\
  \fcol{Tester} & Wajih-Ur-Raza Asif. \\
  \bottomrule
\end{tabular}
\end{center}

\subsection{Black-Box Test Cases}

\subsubsection*{Module: Authentication — Login}
\begin{center}\small
\begin{tabular}{@{}p{0.06\linewidth}p{0.13\linewidth}p{0.30\linewidth}p{0.33\linewidth}p{0.08\linewidth}@{}}
  \toprule
  \textbf{TC \#} & \textbf{Technique} & \textbf{Input} & \textbf{Expected Output} & \textbf{Status} \\
  \midrule
  TC-01 & Valid partition   & \texttt{ali@naawbi.edu} / \texttt{password123}    & Login succeeds; catalog loads; role = instructor  & \textcolor{success}{Pass} \\
TC-02 & Valid partition   & \texttt{ibbi@naawbi.edu} / \texttt{password123}   & Login succeeds; catalog loads; role = student     & \textcolor{success}{Pass} \\
  TC-03 & Invalid partition & \texttt{ali@naawbi.edu} / \texttt{wrongpass}      & Inline error ``Incorrect email or password''      & \textcolor{success}{Pass} \\
TC-04 & Invalid partition & \texttt{nobody@naawbi.edu} / \texttt{password123} & Inline error ``Incorrect email or password''      & \textcolor{success}{Pass} \\
  TC-05 & Boundary / empty  & Empty email + empty password                      & Inline error ``Please fill in all fields''        & \textcolor{success}{Pass} \\
TC-06 & Error guessing    & \texttt{' OR 1=1 --} / any                        & Login fails --- parameterised query blocks SQLi   & \textcolor{success}{Pass} \\
  \bottomrule
\end{tabular}
\end{center}

\subsubsection*{Module: Authentication — Registration}
\begin{center}\small
\begin{tabular}{@{}p{0.06\linewidth}p{0.13\linewidth}p{0.30\linewidth}p{0.33\linewidth}p{0.08\linewidth}@{}}
  \toprule
  \textbf{TC \#} & \textbf{Technique} & \textbf{Input} & \textbf{Expected Output} & \textbf{Status} \\
  \midrule
  TC-07 & Valid partition   & New unique username, email, password, role=Student & Account created; redirected to catalog       & \textcolor{success}{Pass} \\
TC-08 & Invalid partition & Username already taken                            & ``Username or email is already taken''        & \textcolor{success}{Pass} \\
  TC-09 & Invalid partition & Email already taken                                & ``Username or email is already taken''        & \textcolor{success}{Pass} \\
TC-10 & Boundary / empty  & Any field left blank                               & ``Please fill in all fields''                 & \textcolor{success}{Pass} \\
  \bottomrule
\end{tabular}
\end{center}

\subsubsection*{Module: RBAC — Instructor-Only Guard}
\begin{center}\small
\begin{tabular}{@{}p{0.06\linewidth}p{0.13\linewidth}p{0.30\linewidth}p{0.33\linewidth}p{0.08\linewidth}@{}}
  \toprule
  \textbf{TC \#} & \textbf{Technique} & \textbf{Input} & \textbf{Expected Output} & \textbf{Status} \\
  \midrule
  TC-11 & Invalid partition & Student \texttt{ibbi} clicks ``Create Assignment'' & Access Denied alert; no form opens          & \textcolor{success}{Pass} \\
TC-12 & Valid partition   & Instructor \texttt{ali} clicks ``Create Assignment'' & Create Assignment form opens normally       & \textcolor{success}{Pass} \\
  TC-13 & Valid partition   & Student logs in                                  & Catalog shows only enrolled courses (CS101, CS202) & \textcolor{success}{Pass} \\
TC-14 & Valid partition   & Instructor logs in                               & Catalog shows only courses created by \texttt{ali}   & \textcolor{success}{Pass} \\
  \bottomrule
\end{tabular}
\end{center}

\subsubsection*{Module: Grading Interface}
\begin{center}\small
\begin{tabular}{@{}p{0.06\linewidth}p{0.13\linewidth}p{0.30\linewidth}p{0.33\linewidth}p{0.08\linewidth}@{}}
  \toprule
  \textbf{TC \#} & \textbf{Technique} & \textbf{Input} & \textbf{Expected Output} & \textbf{Status} \\
  \midrule
  TC-15 & Valid partition & grade=85 / total=100 / feedback=``Good work'' & Grade saved; submission shows graded in ListView & \textcolor{success}{Pass} \\
TC-16 & Boundary value  & grade = 0                                     & Grade saved as 0; valid                           & \textcolor{success}{Pass} \\
  TC-17 & Boundary value  & grade = total\_points (e.g.\ 100)             & Grade saved at maximum; valid                     & \textcolor{success}{Pass} \\
TC-18 & Boundary value  & grade = total\_points + 1 (e.g.\ 101)         & Inline validation error; Save blocked             & \textcolor{success}{Pass} \\
  TC-19 & Invalid partition & grade = ``abc'' (non-numeric)               & Inline validation error; Save blocked             & \textcolor{success}{Pass} \\
TC-20 & Valid partition & No feedback; grade = 50                     & Grade saved with NULL feedback                    & \textcolor{success}{Pass} \\
  \bottomrule
\end{tabular}
\end{center}

\subsubsection*{Module: Submission Status Derivation}
\begin{center}\small
\begin{tabular}{@{}p{0.06\linewidth}p{0.13\linewidth}p{0.30\linewidth}p{0.33\linewidth}p{0.08\linewidth}@{}}
  \toprule
  \textbf{TC \#} & \textbf{Technique} & \textbf{Input} & \textbf{Expected Output} & \textbf{Status} \\
  \midrule
  TC-21 & Valid partition   & Submission before deadline                        & Badge: \textbf{Submitted}       & \textcolor{success}{Pass} \\
TC-22 & Valid partition   & Submission after deadline; late policy allowed   & Badge: \textbf{Late}            & \textcolor{success}{Pass} \\
  TC-23 & Valid partition   & Deadline passed; no submission; late allowed      & Badge: \textbf{Missing}         & \textcolor{success}{Pass} \\
TC-24 & Valid partition   & Deadline in future; no submission                 & Badge: \textbf{Not Submitted}   & \textcolor{success}{Pass} \\
  TC-25 & Invalid partition & Deadline passed; hard policy; attempt submit      & Submission blocked              & \textcolor{success}{Pass} \\
  \bottomrule
\end{tabular}
\end{center}

\subsection{White-Box Testing}

\subsubsection*{LoginController.handleLogin() --- Branch Coverage}
\begin{quote}\small\ttfamily
B1: email or password empty \hfill $\to$ show error, return\\
B2: User.findByEmail() returns empty \hfill $\to$ show error, return\\
B3: stored hash does not match input hash \hfill $\to$ show error, return\\
B4: all checks pass \hfill $\to$ populate Session, navigate
\end{quote}

\begin{center}
\begin{tabular}{@{}p{0.28\linewidth}p{0.56\linewidth}p{0.10\linewidth}@{}}
  \toprule
  \fcolplain{Branch} & Test Case & Covered \\
  \midrule
  \fcolplain{B1 --- empty fields} & TC-05                  & Yes \\
\fcolplain{B2 --- unknown email} & TC-04                  & Yes \\
  \fcolplain{B3 --- wrong password} & TC-03                  & Yes \\
\fcolplain{B4 --- successful login} & TC-01, TC-02           & Yes \\
  \midrule
  \multicolumn{2}{r}{\textbf{Branch coverage}} & \textbf{4\,/\,4 (100\%)} \\
  \bottomrule
\end{tabular}
\end{center}

\subsubsection*{GradingController.handleSave() --- Branch Coverage}
\begin{quote}\small\ttfamily
B1: grade empty or non-numeric \hfill $\to$ show error, return\\
B2: grade $>$ total\_points \hfill $\to$ show error, return\\
B3: grade $<$ 0 \hfill $\to$ show error, return\\
B4: all checks pass \hfill $\to$ Submission.saveGrade(), refresh
\end{quote}

\begin{center}
\begin{tabular}{@{}p{0.28\linewidth}p{0.56\linewidth}p{0.10\linewidth}@{}}
  \toprule
  \fcolplain{Branch} & Test Case & Covered \\
  \midrule
  \fcolplain{B1 --- non-numeric input} & TC-19                        & Yes \\
\fcolplain{B2 --- grade over maximum} & TC-18                        & Yes \\
  \fcolplain{B3 --- grade below minimum} & Boundary: grade = $-1$ (manual) & Yes \\
\fcolplain{B4 --- valid save} & TC-15, TC-16, TC-17          & Yes \\
  \midrule
  \multicolumn{2}{r}{\textbf{Branch coverage}} & \textbf{4\,/\,4 (100\%)} \\
  \bottomrule
\end{tabular}
\end{center}

\subsubsection*{Assignment.fetchWithStatusForUser() --- Statement Coverage}
The method executes a LEFT JOIN and derives submission status via conditional
logic. All four status paths (Submitted, Late, Missing, Not Submitted) are
exercised by TC-21 through TC-24 against the seeded dataset --- achieving
statement coverage of every reachable branch in the status derivation block.

\subsection{Bug Report}

\begin{center}\small
\begin{tabular}{@{}p{0.07\linewidth}p{0.11\linewidth}p{0.58\linewidth}p{0.06\linewidth}p{0.08\linewidth}@{}}
  \toprule
  \textbf{Bug ID} & \textbf{Module} & \textbf{Description and Fix} & \textbf{Sev.} & \textbf{Status} \\
  \midrule
  BUG-01 & Grading     & \texttt{GradingView} threw NPE on an assignment with zero submissions. \emph{Fix:} empty-list guard before ListView population.                                            & P1 & \textcolor{success}{Fixed} \\
  BUG-02 & RBAC        & Instructors saw courses they did not create because \texttt{created\_by} was not being set. \emph{Fix:} \texttt{CreateCourseController} now sets it from \texttt{Session}.     & P1 & \textcolor{success}{Fixed} \\
  BUG-03 & Gradebook   & GridPane column headers misaligned when a student had no submissions for the first assignment. \emph{Fix:} iterate all assignments independent of submission presence.        & P2 & \textcolor{success}{Fixed} \\
  BUG-04 & Attachments & Attachment chip did not render after stream refresh. \emph{Fix:} stream loader now calls \texttt{fetchByAnnouncementId()} per announcement. & P1 & \textcolor{success}{Fixed} \\
  \bottomrule
\end{tabular}
\end{center}

\subsection{Test Execution Results}

Testing was performed manually against the seeded database (\texttt{make seed})
using all three test accounts.

\begin{center}
\small\begin{tabular}{@{}p{0.24\linewidth}p{0.12\linewidth}p{0.58\linewidth}@{}}
  \toprule
  \fcolplain{Account} & Role       & Scenarios Tested \\
  \midrule
  \texttt{ali@naawbi.edu}  & Instructor & Login, RBAC guard, create assignment, post announcement with both attachment types, grade all 6 submissions, view class gradebook, People tab, logout. \\
  \texttt{ibbi@naawbi.edu} & Student    & Login, register, enrolment filtering, submit assignment, view submission status, personal gradebook, To-Do list, logout. \\
  \texttt{sara@naawbi.edu} & Student    & Login, second-student perspective, verify gradebook isolation, People tab roster. \\
  \bottomrule
\end{tabular}
\end{center}

\noindent\textbf{Summary.} All 25 black-box test cases passed. All 4 identified
bugs (BUG-01 through BUG-04) were fixed before submission. No unfixed P1 bugs
remain. One P2 issue (BUG-03) was found and fixed. System behaviour matched
acceptance criteria across all Sprint~3 user stories.

\newpage
% ═══════════════════════════════════════════════════════════════
%                 PART C — DELIVERABLE 3 DOCUMENT
% ═══════════════════════════════════════════════════════════════
\begin{center}
  {\LARGE\bfseries\color{primary} PART C\,\textemdash\,Deliverable 3 Document}\\[4pt]
  {\color{muted}\small Full system documentation after the final sprint}\\[4pt]
  {\color{rule!40}\rule{7cm}{0.6pt}}
\end{center}

% ───────────────────────────────────────────────────────────────
\section{Introduction}

\subsection{Context}

Naawbi is a university-level Learning Management System built iteratively by
Team Ibwaan for CS3009 Software Engineering (Spring~2026). This document covers
Deliverable~3, the output of Sprint~3 and the final sprint of the project.

\subsection{Progress from Sprint 2 to Sprint 3}

Sprint~2 delivered the content and submission layer --- announcements,
assignments, file-based submissions, submission status tracking, and the
student To-Do list. All of that ran without any authentication; any user could
reach any feature. Sprint~3 closes the platform: every user now registers with
an identity and role, logs in to a session-guarded workspace, and participates
in a complete grade cycle --- from submission through instructor grading to a
personal gradebook --- while role-based access control prevents students from
accessing instructor features and vice versa.

Sprint~3 also completed two features that were implied by D2's assignment
workflow but were never formally built: the Assignments tab (giving students a
per-course view of all assignments with submission status) and the People tab
(showing the course roster with role badges).

\subsection{Carried from Sprint 2 Backlog}

\begin{center}
\begin{tabular}{@{}p{0.14\linewidth}p{0.78\linewidth}@{}}
  \toprule
  \textbf{User Story} & \textbf{Description} \\
  \midrule
  US\,1--4   & Authentication --- register, login, role assignment, logout \\
US\,12--13 & Announcement attachments --- file and link \\
  US\,19--22 & Grading interface and gradebook (instructor + student views) \\
  \bottomrule
\end{tabular}
\end{center}

\noindent Sprint~3 additionally scoped US\,15--16 (Assignments tab + People
tab), which had not appeared in any prior backlog.

% ───────────────────────────────────────────────────────────────
\section{Project Vision (Refined)}

Build a clean, reliable, course-centric LMS that instructors use to manage
content and evaluations, and that students use to track progress and
performance --- without needing external tools for any core academic workflow.

\medskip
\noindent Sprint~1 built the course structure. Sprint~2 populated it with
content and submissions. Sprint~3 closes the loop: authenticated identities,
enforced roles, grading, and a gradebook. \textbf{The full feature set is now
delivered.}

% ───────────────────────────────────────────────────────────────
\section{Intended Use of the System}

\subsection{Stakeholders and Sprint 3 Impact}

\begin{center}
\begin{tabular}{@{}p{0.18\linewidth}p{0.13\linewidth}p{0.62\linewidth}@{}}
  \toprule
  \fcol{Stakeholder} & Primary Need & Sprint 3 Impact \\
  \midrule
  \fcol{Student} & Single place to track work, deadlines, and grades  & Logs in with a verified account; sees only enrolled courses; submits assignments; views per-course assignment status in the Assignments tab; sees personal grades and feedback in the Gradebook. \\
\fcol{Instructor} & Unified workflow for content, collection, and evaluation & Logs in; is the only role that can post, create assignments, and grade; can view a full class-wide gradebook at a glance; People tab shows course roster. \\
  \fcol{System} & Consistent enforcement of access rules             & RBAC enforced at every protected action via \texttt{Session} role check; no unauthenticated code path remains. \\
  \bottomrule
\end{tabular}
\end{center}

\noindent No stakeholder needs changed during Sprint~3; the existing three
roles map cleanly onto the newly-implemented authentication and RBAC.

% ───────────────────────────────────────────────────────────────
\section{Features and Overall Functionality}

\subsection{Cumulative Feature Summary (Deliverables 1--3)}

Sprint~3 delivered the final third of the system. The table below shows every
feature across all three deliverables and which sprint produced it.

\begin{center}\small
\begin{tabular}{@{}p{0.48\linewidth}p{0.22\linewidth}p{0.12\linewidth}p{0.10\linewidth}@{}}
  \toprule
  \textbf{Feature} & \textbf{User Stories} & \textbf{Deliv.} & \textbf{Status} \\
  \midrule
  Course Creation                             & US\,5 (D1)              & D1 & \textcolor{success}{Done} \\
Course Catalogue and Home Page      & US\,6, US\,9            & D1 & \textcolor{success}{Done} \\
  Unique Class Code Enrollment                & US\,7, US\,8            & D1 & \textcolor{success}{Done} \\
Course Announcements                & US\,10, US\,11          & D2 & \textcolor{success}{Done} \\
  Assignment Creation with Policies           & US\,14, US\,15 (D2)     & D2 & \textcolor{success}{Done} \\
File Submissions and Status Tracking & US\,16, US\,17, US\,18 & D2 & \textcolor{success}{Done} \\
  Threaded Comments on Submissions            & US\,16 (extension)      & D2 & \textcolor{success}{Done} \\
To-Do List and Deadline Awareness   & US\,23, US\,24          & D2 & \textcolor{success}{Done} \\
  User Registration and Login                 & US\,1, US\,2, US\,3, US\,4 & D3 & \textcolor{success}{Done} \\
Logout                              & US\,4                   & D3 & \textcolor{success}{Done} \\
  Role-Based Access Control                   & US\,5 (D3)              & D3 & \textcolor{success}{Done} \\
Announcement Attachments (File + Link) & US\,12, US\,13       & D3 & \textcolor{success}{Done} \\
  Assignments Tab with Status Badges          & US\,15 (D3)             & D3 & \textcolor{success}{Done} \\
People Tab (Course Roster)          & US\,16 (D3)             & D3 & \textcolor{success}{Done} \\
  Instructor Grading Interface                & US\,19, US\,20          & D3 & \textcolor{success}{Done} \\
Student Personal Gradebook          & US\,21                  & D3 & \textcolor{success}{Done} \\
  Instructor Class Gradebook                  & US\,22                  & D3 & \textcolor{success}{Done} \\
  \bottomrule
\end{tabular}
\end{center}

\subsection{How Sprint 3 Enhanced the Final Integrated Product}

\begin{itemize}
  \item \textbf{Every Sprint~1 and Sprint~2 feature is now guarded.}
        Authentication and RBAC retrofit protects every earlier controller
        --- course creation, assignment creation, submission, and To-Do list
        are all now session-bound.
  \item \textbf{Complete grade cycle.} A student can now go from enrolment
        through submission to a graded result with instructor feedback ---
        entirely within Naawbi. No external tool is required.
  \item \textbf{Instructor macro-view.} The class-wide gradebook and People
        tab give instructors visibility over student population and
        performance, which the earlier sprints lacked.
  \item \textbf{Richer content.} File and link attachments on announcements
        close a gap from D2 where instructors could post text but not share
        documents or external references inline.
  \item \textbf{Accessible workload view.} The Assignments tab consolidates
        assignment status into one glanceable place per course, complementing
        the global To-Do list from D2.
\end{itemize}

% ───────────────────────────────────────────────────────────────
\section{User Stories (Sprint 3 + Carried from Sprint 2)}
\label{sec:sprint3-stories}

Sprint~3 completed both newly-scoped stories and Sprint~2 backlog
carry-forward. Detailed backlog tasks and structured specifications appear in
Part~B.

\begin{itemize}
  \item \textbf{Carried from Sprint 2:} US\,1--4 (Authentication --- register, login, role assignment, logout), US\,12--13 (Announcement attachments --- file and link), US\,19--22 (Grading interface, student gradebook, instructor gradebook).
  \item \textbf{New in Sprint 3:} US\,5 (Role-Based Access Control), US\,15 (Assignments tab with status badges), US\,16 (People tab with course roster).
\end{itemize}

\noindent All ten stories are written in standard \emph{``As a \ldots, I want
\ldots, so that \ldots''} form in Part~B, Section~``User Stories Selected for
Sprint~3''. No story from any earlier sprint remains open.

% ───────────────────────────────────────────────────────────────
\section{Non-Functional Requirements (Final State)}

\begin{center}
\begin{tabular}{@{}p{0.22\linewidth}@{\hspace{6pt}}p{0.72\linewidth}@{}}
  \toprule
  \fcol{NFR Category} & Requirement \\
  \midrule
  Security
    & Passwords are never stored or compared in plain text ---
      SHA-256 hashing is applied at all registration and login points via
      \texttt{PasswordUtil.hash()}. Session state is held in a singleton that
      is explicitly cleared on logout. No credentials persist beyond the
      session. RBAC guards are applied at the controller level before every
      instructor-only action. \\
Performance
    & Gradebook queries (class-wide grid) execute in a single JOIN
      (\texttt{fetchGradebookForCourse}) rather than N+1 queries per student.
      All primary views target loading within 2--3 seconds on typical lab
      hardware at class sizes up to $\sim$200 students per course. \\
  Reliability
    & Grade saves are atomic --- \texttt{Submission.saveGrade()} issues a
      single parameterised UPDATE; no partial grade state is possible. Session
      logout completes synchronously before the \texttt{LoginView} is
      displayed. Submission timestamps are set server-side, not client-side,
      ensuring consistent status evaluation. \\
Usability
    & Role-adaptive UI: instructor-only buttons (Post Announcement, Create
      Assignment, Grade Submissions) are conditionally rendered --- not merely
      disabled --- for students. Status badges (Submitted / Late / Missing /
      Not Submitted) use both colour and text so they are never colour-only.
      Gradebook cells show ``---'' for ungraded rather than blank or null. \\
  Maintainability
    & All Sprint~3 features follow the established layered MVC architecture
      (controller / model / view). RBAC checks, grade validation, and gradebook
      queries live in model/controller classes, not in FXML event handlers. All
      new views follow the same FXML + CSS structure as existing ones. \\
Scalability
    & Supports multiple courses and concurrent users typical of a
      department-scale deployment (not enterprise scale). The JOIN-based
      gradebook query avoids per-student round-trips that would not scale past
      a single classroom. \\
  Auditability
    & Graded submissions retain the original submission timestamp alongside
      the new grade and feedback fields --- no information is overwritten.
      Session login is gated on credential verification; no anonymous access
      path remains in the application after Sprint~3. \\
  \bottomrule
\end{tabular}
\end{center}

% ───────────────────────────────────────────────────────────────
\section{Testing and Quality Assurance}

Full test artefacts (test plan table, 25 test cases, branch-coverage tables,
and bug report table) are in Part~B, Section~\ref{sec:testing}. This section
summarises each of the six required items.

\subsection{Test Plan}
Scope: all Sprint~3 features (Authentication, RBAC, Announcement Attachments,
Assignments Tab, People Tab, Grading, Gradebook). Environment: Windows~11,
Java~24, JavaFX~21, PostgreSQL~16 seeded via \texttt{data/seed.sql}. Entry
criteria: \texttt{make seed} and \texttt{make compile} succeed. Exit criteria:
all P1 test cases pass; no unfixed P1 bugs.

\subsection{Black-Box Test Cases}
25 test cases were designed using equivalence partitioning, boundary value
analysis, and error guessing across four modules (Login, Registration,
RBAC, Grading, Submission Status). All 25 passed. Coverage highlights:
valid/invalid credentials, empty fields, SQL injection (TC-06), grade at
boundaries ($0$, \texttt{total\_points}, \texttt{total\_points}+1), and all
four submission status states.

\subsection{White-Box Testing Evidence}
Branch coverage of \texttt{LoginController.handleLogin()} = 4/4 (100\%) and
\texttt{GradingController.handleSave()} = 4/4 (100\%). Statement coverage of
\texttt{Assignment.fetchWithStatusForUser()} exercises all four status
derivation paths via TC-21 through TC-24.

\subsection{Bug Report Summary}

\begin{center}\small
\begin{tabular}{@{}p{0.08\linewidth}p{0.13\linewidth}p{0.57\linewidth}p{0.12\linewidth}@{}}
  \toprule
  \textbf{Bug ID} & \textbf{Module} & \textbf{Description (short)} & \textbf{Status} \\
  \midrule
  BUG-01 & Grading     & NPE on empty submission list                          & \textcolor{success}{Fixed} \\
  BUG-02 & RBAC        & Instructor saw courses they did not create           & \textcolor{success}{Fixed} \\
  BUG-03 & Gradebook   & GridPane column misalignment with missing subs.      & \textcolor{success}{Fixed} \\
  BUG-04 & Attachments & Attachment chips not rendering after stream refresh  & \textcolor{success}{Fixed} \\
  \bottomrule
\end{tabular}
\end{center}

\subsection{Test Results and Observations}
All 25 black-box cases passed; both target methods hit 100\% branch coverage.
The SQL injection test (TC-06) confirmed that parameterised queries block
injection at the login boundary. Boundary tests (TC-16--18) confirmed grade
$=0$ and grade $=$ \texttt{total\_points} are accepted while
\texttt{total\_points}$+1$ is rejected.

\subsection{How Testing Improved System Quality}
The four bugs caught in test would all have been user-facing in submission:
BUG-01 would have crashed the Grading modal for any assignment with zero
submissions (a common early state); BUG-02 broke the core RBAC invariant
entirely; BUG-04 made the attachment feature invisible to students. Explicit
coverage of all four submission-status paths also surfaced a gap in the
``Not Submitted'' derivation that was corrected before integration.

% ───────────────────────────────────────────────────────────────
\section{Iteration 3 Implementation Screenshots}

\begin{figure}[H]
  \centering
  \begin{minipage}{0.48\textwidth}\centering
    \imgOrPlaceholder[width=\textwidth]{assets/login_view.png}{Login View}
    \subcaption{LoginView --- tabbed Sign In / Register.}
  \end{minipage}\hfill
  \begin{minipage}{0.48\textwidth}\centering
    \imgOrPlaceholder[width=\textwidth]{assets/course_catalog_role.png}{Course Catalog}
    \subcaption{Course Catalog filtered by role (instructor).}
  \end{minipage}
  \caption{Authentication and role-based catalog filtering.}
\end{figure}

\begin{figure}[H]
  \centering
  \imgOrPlaceholder[width=0.72\textwidth]{assets/announcement_attachments.png}{Announcement Attachments}
  \caption{\textbf{Course Stream} --- announcement with file and link attachment chips stored in \texttt{announcement\_attachments}.}
\end{figure}

\begin{figure}[H]
  \centering
  \begin{minipage}{0.48\textwidth}\centering
    \imgOrPlaceholder[width=\textwidth]{assets/assignments_tab_student.png}{Assignments Tab — Student}
    \subcaption{Student view with status badges.}
  \end{minipage}\hfill
  \begin{minipage}{0.48\textwidth}\centering
    \imgOrPlaceholder[width=\textwidth]{assets/assignments_tab_instructor.png}{Assignments Tab — Instructor}
    \subcaption{Instructor view with Grade Submissions buttons.}
  \end{minipage}
  \caption{Assignments Tab --- role-adaptive cards sharing the same FXML backbone.}
\end{figure}

\begin{figure}[H]
  \centering
  \imgOrPlaceholder[width=0.72\textwidth]{assets/grading_modal.png}{Grading Modal}
  \caption{\textbf{GradingView} --- ListView of submissions on the left; numeric grade and feedback editor on the right.}
\end{figure}

\begin{figure}[H]
  \centering
  \begin{minipage}{0.48\textwidth}\centering
    \imgOrPlaceholder[width=\textwidth]{assets/student_gradebook.png}{Student Gradebook}
    \subcaption{Student personal gradebook (read-only).}
  \end{minipage}\hfill
  \begin{minipage}{0.48\textwidth}\centering
    \imgOrPlaceholder[width=\textwidth]{assets/instructor_gradebook.png}{Instructor Class Gradebook}
    \subcaption{Instructor class-wide grid from \texttt{fetchGradebookForCourse()}.}
  \end{minipage}
  \caption{Gradebook views --- role-adaptive rendering (table vs.\ GridPane).}
\end{figure}

\begin{figure}[H]
  \centering
  \imgOrPlaceholder[width=0.72\textwidth]{assets/people_tab.png}{People Tab}
  \caption{\textbf{People Tab} --- role-grouped roster cards from \texttt{Course.fetchEnrolledUsers()}.}
\end{figure}

% ───────────────────────────────────────────────────────────────
\section{Work Division for Sprint 3}

\begin{center}
\begin{tabular}{@{}p{0.22\linewidth}@{\hspace{6pt}}p{0.72\linewidth}@{}}
  \toprule
  \fcolplain{Member} & Sprint 3 Responsibilities \\
  \midrule
  Ibraheem Farooq
    & Sprint~3 planning and GitHub Projects board management across all sprint
      stages. Branch and PR oversight --- reviewed and merged
      \texttt{feature/auth}, \texttt{feature/enrollment-rbac}, and
      \texttt{feature/announcements-materials} PRs. Implemented end-to-end
      authentication and RBAC: tabbed Login/Register UI, \texttt{Session}
      singleton, \texttt{PasswordUtil}, \texttt{LoginController},
      enrollment-based catalog filtering, and Access Denied guard on
      instructor-only actions. Implemented announcement file and link
      attachments (US\,12 \& 13): \texttt{CreateAnnouncement} extension,
      \texttt{announcement\_attachments} table, and attachment chip rendering
      in the course stream. Final review and packaging of deliverable
      artefacts. \\
  \addlinespace[4pt]
Wajih-Ur-Raza Asif
    & Sprint~3 user stories, sub-story breakdowns, and all six structured
      specifications. Consistency verification between acceptance criteria and
      actual implemented behaviour. Deliverable~3 document authoring and
      completeness review. Implemented the Assignments tab and People tab
      (US\,15 \& 16): \texttt{Assignment.fetchWithStatusForUser()}, status
      badge rendering, Grade Submissions button,
      \texttt{Course.fetchEnrolledUsers()}, and roster card rendering.
      Implemented the full grading pipeline (US\,19--22):
      \texttt{grade}/\texttt{feedback} DB schema, \texttt{Submission} model
      extensions (\texttt{fetchByAssignmentId}, \texttt{saveGrade},
      \texttt{fetchGradedForStudent}, \texttt{fetchGradebookForCourse}),
      \texttt{GradingView.fxml} + \texttt{GradingController}, and
      \texttt{GradebookView.fxml} + \texttt{GradebookController} (student
      personal view and instructor class-wide grid). End-to-end test plan
      execution using seeded test data. \\
  \addlinespace[4pt]
  Ali Aan Khowaja
    & Lead implementation of all Sprint~2 features: Announcements, Assignment
      Creation, AssignmentDetails with file submission flow and threaded
      comments, and the To-Do List. UI design and screen layout for all
      Sprint~2 and Sprint~3 views. Produced Sprint~3 implementation
      screenshots. \\
  \bottomrule
\end{tabular}
\end{center}

% ───────────────────────────────────────────────────────────────
\section{Retrospective}

\textbf{What went well.} All Sprint~3 stories shipped on time, including the
two additions (Assignments tab, People tab). The layered MVC architecture
established in Sprint~1 absorbed every new feature without structural change.
The idempotent \texttt{seed.sql} with three test accounts and mixed-state
submissions accelerated manual testing, and the branch-per-feature strategy
kept code review manageable.

\textbf{What could be improved.} Authentication should have been Sprint~1
work, not Sprint~3 --- RBAC had to be retrofitted onto every existing
controller. The grading chain (db $\rightarrow$ ui $\rightarrow$ gradebook)
serialised work that could have run in parallel, and US\,15--16 (Assignments
and People tabs) should have been identified during D2 planning when the
assignment infrastructure was fresh.

\textbf{What we would do differently.} Ship authentication in Sprint~1,
identify course-home navigation during D2, and use an integration branch for
feature chains with linear merge dependencies.

% ───────────────────────────────────────────────────────────────
\section{Final System Summary}

Naawbi LMS was built across three sprints, each delivering a distinct layer
of the final product.

\begin{center}
\begin{tabular}{@{}p{0.18\linewidth}p{0.13\linewidth}p{0.62\linewidth}@{}}
  \toprule
  \fcol{Sprint} & Deliverable & Core Contribution \\
  \midrule
  \fcol{Sprint 1} & D1 & Course structure --- creation, catalogue, class-code enrolment. \\
  \fcol{Sprint 2} & D2 & Content and submissions --- announcements, assignments, file uploads, status tracking, To-Do list. \\
  \fcol{Sprint 3} & D3 & Identity and evaluation --- authentication, RBAC, attachments, Assignments/People tabs, grading, gradebook. \\
  \bottomrule
\end{tabular}
\end{center}

\noindent The final system fulfils its original vision: a course-centric LMS
where instructors manage content and evaluations, and students track progress
and performance without needing any external tool. \textbf{All planned user
stories across all three sprints have been delivered and verified.}

\end{document}
